Residual diagnostics show that the declared fitted forms are not equally
compatible with the reconstructed line. Values below 1 in the table
indicate that the RMS residual is no larger than the local reconstructed
line half-width on average.

\begin{longtable}[]{@{}lrr@{}}
\toprule\noalign{}
Fitted form & MZ226 normalized RMS & MZ310 normalized RMS \\
\midrule\noalign{}
\endhead
\bottomrule\noalign{}
\endlastfoot
Free exponent & 0.69 & 0.75 \\
\(p=0.5\) & 5.67 (boundary-limited) & 4.88 \\
\(p=0.7\) & 2.47 & 2.62 \\
Source-panel exponent & 0.82 & 1.50 \\
\(p=1\) & 0.78 & 0.84 \\
Chu Kane-region & 6.08 (boundary-limited) & 7.66 (boundary-limited) \\
\end{longtable}

Energy-resolved residuals reveal where the fitted forms fail rather than
compressing model adequacy into one scalar. The \(p=0.5\), \(p=0.7\),
and Chu candidates show substantial structured residuals; the
free-exponent and \(p=1\) forms remain much closer to the reconstructed
line over the nominal domain.

\begin{figure}
\centering
\includegraphics{figures/figure6_residuals_mz226.pdf}
\caption{MZ226 residuals versus photon energy. The shaded band is the
local reconstructed line half-width around zero residual.}
\end{figure}

\begin{figure}
\centering
\includegraphics{figures/figure7_residuals_mz310.pdf}
\caption{MZ310 residuals versus photon energy. The shaded band is the
local reconstructed line half-width around zero residual.}
\end{figure}

\FloatBarrier

The nominal line-envelope screen retains the free-exponent,
source-panel, and \(p=1\) forms for MZ226, giving 2.61 meV, and the
free-exponent and \(p=1\) forms for MZ310, giving 2.09 meV. The
threshold grid shows why these values must remain secondary.

\begin{longtable}[]{@{}
  >{\raggedright\arraybackslash}p{(\columnwidth - 4\tabcolsep) * \real{0.2727}}
  >{\raggedleft\arraybackslash}p{(\columnwidth - 4\tabcolsep) * \real{0.3636}}
  >{\raggedleft\arraybackslash}p{(\columnwidth - 4\tabcolsep) * \real{0.3636}}@{}}
\toprule\noalign{}
\begin{minipage}[b]{\linewidth}\raggedright
Screen rule
\end{minipage} & \begin{minipage}[b]{\linewidth}\raggedleft
MZ226 retained span
\end{minipage} & \begin{minipage}[b]{\linewidth}\raggedleft
MZ310 retained span
\end{minipage} \\
\midrule\noalign{}
\endhead
\bottomrule\noalign{}
\endlastfoot
RMS \(\leq0.75\); coverage \(\geq0.90\)-\(0.975\) & one model; no span &
no model retained \\
RMS \(\leq1.0\)-\(1.5\); coverage \(\geq0.90\) or \(0.95\) & 2.61 meV &
2.09 meV \\
RMS \(\leq1.0\)-\(1.5\); coverage \(\geq0.975\) & 0.63 meV & 2.09 meV \\
\end{longtable}

The screen is therefore informative about grossly incompatible models
but its numerical span is not invariant to the descriptive thresholds.
It is excluded from the abstract-level headline and is not interpreted
as a preferred physical interval.

\begin{figure}
\centering
\includegraphics{figures/figure3_relative_fitted_intercepts.pdf}
\caption{Nominal fitted-intercept positions relative to the lowest
non-boundary model for each specimen.}
\end{figure}

\begin{figure}
\centering
\includegraphics{figures/figure5_model_residual_compatibility.pdf}
\caption{RMS residuals normalized by the local reconstructed line
half-width.}
\end{figure}

The robust conclusion is therefore not that one universal 5-7 meV model
uncertainty applies to HgCdTe. It is that the reported intercept is
materially conditioned by model form and especially by fit-domain and
weighting choices, while the isolated 0.198 eV reconstruction
irregularity is not the controlling factor.

\subsection{\texorpdfstring{5.4 Fixed-\(\alpha\) operational
coordinates}{5.4 Fixed-\textbackslash alpha operational coordinates}}\label{fixed-alpha-operational-coordinates}

Fixed-absorption crossings vary substantially with the selected
threshold, as expected because they are different operational
coordinates. They are useful only when the threshold is stated and held
constant. They are not included in the fitted-model span and are not
used to rank Hansen, Seiler, Laurenti, or any provisional relation.

\subsection{5.5 Comparison scale, not equation
ranking}\label{comparison-scale-not-equation-ranking}

At the nominal reported compositions and 300 K, the Hansen-Seiler
separations are approximately 0.255 meV for MZ226 and 0.177 meV for
MZ310. These values are smaller than:

\begin{itemize}
\tightlist
\item
  the boundary-excluded fitted-model spans;
\item
  the coherent calibration-corner sensitivity;
\item
  the composition-energy scale associated with half of the last reported
  composition digit.
\end{itemize}

This comparison demonstrates that the available records do not support a
sub-meV equation ranking. It does not establish physical equivalence of
the equations.
